\section{Giới thiệu}

Một hệ thống đa tác tử (Multi-Agent System -- MAS) dựa trên mô hình ngôn ngữ lớn (LLM) chia một nhiệm vụ khó cho nhiều tác tử, mỗi tác tử giữ một vai trò riêng và trao đổi kết quả trung gian với nhau. Cách làm này đạt hiệu quả cao trên nhiều tác vụ suy luận phức tạp, nhưng nó tạo ra một vấn đề đi kèm: các tác tử nói nhiều với nhau, và những gì chúng nói tích tụ thành ngữ cảnh rất dài. Khi lịch sử trao đổi dài ra, câu trả lời của một tác tử có thể phụ thuộc vào vài câu nhỏ nằm giữa hàng nghìn từ nhiễu. Hiện tượng mô hình bỏ sót thông tin nằm giữa ngữ cảnh dài đã được mô tả trong tài liệu với tên gọi ``lost in the middle''~\cite{liu2024lost}.

Hai bài báo được tìm hiểu trong báo cáo này tiếp cận vấn đề trên từ hai phía khác nhau. \textbf{Agent-Radar}~\cite{agentradar2026} xử lý phía \emph{đọc}: khi một tác tử chuẩn bị sinh câu trả lời, lịch sử trao đổi nên được đọc theo cách có chọn lọc thay vì đọc đều như nhau. Phương pháp này chấm điểm từng câu theo độ liên quan ngữ nghĩa, điều chỉnh theo khoảng cách giữa các tác tử trên đồ thị giao tiếp và theo tuổi của thông điệp, rồi lái attention của LLM về các câu được chọn. \textbf{DAIN}~\cite{dain2026} xử lý phía \emph{tham gia}: trong bài toán hợp nhất dữ liệu đa phương thức, không phải mọi chuyên gia đều cần chạy cho mọi mẫu đầu vào. DAIN dùng một bộ điều phối học được để kích hoạt thưa một nhóm tác tử chuyên biệt cho từng mẫu và điều phối trao đổi thông tin nén giữa chúng.

Hai công trình này được chọn vì chúng bổ sung cho nhau theo một hướng khá tự nhiên: một bên quản lý \emph{thông tin} luân chuyển giữa các tác tử, một bên quản lý \emph{thành phần} tham gia suy luận. Cả hai đều đưa ra cách chọn lọc phụ thuộc ngữ cảnh thay vì chính sách tĩnh, và cả hai đều có kết quả thực nghiệm chi tiết để đối chiếu.

Phạm vi của báo cáo giới hạn ở việc tìm hiểu và minh họa hai công trình nêu trên; đây là bài báo cáo học tập, không phải công trình đề xuất thuật toán mới. Cụ thể, báo cáo gồm các phần sau. Phần~\ref{sec:background} nhắc lại các khái niệm cần thiết. Phần~\ref{sec:agentradar} và phần~\ref{sec:dain} phân tích từng phương pháp. Phần~\ref{sec:comparison} so sánh hai thiết kế. Phần~\ref{sec:experiments} mô tả hai chương trình thực nghiệm, và phần~\ref{sec:results} trình bày kết quả đo được. Phần~\ref{sec:conclusion} kết luận.
